\documentclass[12pt,a4paper]{article}

\usepackage[utf8]{inputenc}
\usepackage[T1]{fontenc}
\usepackage{amsmath,amssymb}
\usepackage{booktabs}
\usepackage{hyperref}
\usepackage[margin=1in]{geometry}
\usepackage{natbib}
\usepackage{xcolor}
\usepackage{array}
\usepackage{caption}

\hypersetup{
    colorlinks=true,
    linkcolor=blue!60!black,
    citecolor=green!50!black,
    urlcolor=blue!70!black,
}

\title{Decentralized Cooperative Positioning via Uncertainty-Bounded Fusion:\\
A Prototype Toward Sensor-Agnostic Navigation Beyond GNSS}

\author{
Tristan Stoltz\textsuperscript{1}\\[4pt]
\textsuperscript{1}Luminous Dynamics\\
\texttt{tristan.stoltz@evolvingresonantcocreationism.com}
}

\date{April 2026}

\begin{document}
\maketitle

\begin{abstract}
We present a decentralized cooperative positioning prototype designed to preserve short-duration location continuity when GNSS is degraded or denied, using commodity-device sensors (BLE, WiFi RTT, UWB, accelerometer, gyroscope, barometer) and peer-to-peer ranging between cooperating devices. The current implementation is best understood as a foreground-first temporal bridge for indoor transitions, underground dead zones, and other localized outages rather than as an indefinite infrastructure-free replacement for GNSS.

The architecture consists of three layers: (1)~a pure-Rust positioning library implementing Gauss-Newton weighted trilateration, uncertainty-aware ranging models, a local Extended Kalman Filter, covariance-intersection helpers for decentralized fusion, and bounded Pedestrian Dead Reckoning (PDR); (2)~a Holochain-based registry for admission, provenance, and delayed audit of anchors and measurements; and (3)~a mobile integration bridge supporting multiple terrestrial ranging technologies. The same mathematical core is expressed through a sensor-agnostic measurement model so that future deployments can replace phone radios with optical links, star trackers, terrain-relative navigation, and celestial timing.

Benchmarks on release-compiled Rust show 3D trilateration at 11.5~\textmu s/call and PDR propagation at 0.3~\textmu s/step, demonstrating that the computational core is not the bottleneck on commodity hardware. These are algorithmic benchmarks only; real-world accuracy remains governed by radio propagation, operating-system policy, and sensor physics.

The key contribution is not a claim of solved decentralized trust, but a separation of concerns: fast local estimation, explicit uncertainty models, conservative peer-fusion primitives, and a decentralized provenance layer. This creates a credible path from a terrestrial smartphone prototype to a sensor-agnostic cooperative navigation kernel for lunar, Martian, and deep-space environments where GNSS does not exist.
\end{abstract}

\textbf{Keywords:} decentralized positioning, cooperative localization, GNSS-denied navigation, covariance intersection, pedestrian dead reckoning, Holochain, BLE Channel Sounding, space navigation

\section{Introduction}

Global Navigation Satellite Systems (GNSS) are a single point of failure embedded in nearly every layer of modern infrastructure. In May 2024, a solar storm disrupted GPS-based tractor guidance systems across the American Midwest, halting spring planting overnight. The U.S. Department of Homeland Security estimates that a nationwide GPS outage would cost the economy over \$1 billion per day \citep{dhs2021pnt}. A Carrington-class coronal mass ejection---an event with estimated recurrence of 100--200 years---would degrade or destroy the entire GNSS constellation for hours to days \citep{cannon2013extreme}.

Existing alternatives to GNSS (eLoran, Iridium STL, NextNav) share a critical limitation: they require centralized infrastructure controlled by a single entity. When that entity fails, the positioning service fails with it.

We propose and implement a different approach: \emph{decentralized cooperative positioning}, where participating devices collectively constrain position through peer-to-peer ranging, bounded measurement fusion, and distributed anchor management. The current contribution is a prototype architecture and computational core, not a claim that today's commodity phones can autonomously sustain an indefinite public positioning mesh under all operating-system and RF constraints.

The system addresses three scenarios:
\begin{enumerate}
    \item \textbf{GPS augmentation}: When GNSS is available, phones with GPS act as anchors for phones without (indoor users, underground, urban canyons).
    \item \textbf{GPS denial}: When GNSS is lost (solar storm, jamming, spoofing), nearby cooperative measurements and bounded inertial propagation can preserve continuity for a short operational window.
    \item \textbf{No-GNSS environments}: Lunar, Martian, and deep-space systems where the same mathematical core can be re-used with different sensors: optical links, star trackers, celestial timing, and terrain-relative navigation.
\end{enumerate}

\subsection{Contributions}

\begin{enumerate}
    \item A pure-Rust positioning library implementing body-agnostic trilateration, adaptive PDR, multiple terrestrial and non-terrestrial ranging models, a sensor-agnostic measurement abstraction, and covariance-intersection helpers
    \item A decentralized registry architecture in which Holochain carries admission, provenance, and delayed audit metadata rather than serving as the real-time fusion engine itself
    \item A mobile integration bridge with 5 positioning modes (GPS, Hybrid, Cooperative Mesh, Dead Reckoning, Offline) appropriate to a foreground-first terrestrial prototype
    \item Benchmarks demonstrating that the computational core is lightweight enough for mobile and embedded deployment: 11.5~\textmu s trilateration and 0.3~\textmu s PDR step
\end{enumerate}

\section{Related Work}

\paragraph{Cooperative positioning.} Wymeersch et al.\ \citep{wymeersch2009cooperative} established the theoretical framework for cooperative localization, showing that peer-to-peer ranging can improve positioning accuracy by 74\% over infrastructure-only approaches. Recent work by the ISPRS \citep{isprs2025uwb} demonstrates UWB/WiFi RTT hybrid cooperative positioning with Split Covariance Intersection for decentralized fusion. This distinction matters: trust or reputation can influence admission and measurement priors, but it is not a mathematical substitute for correlation-consistent fusion.

\paragraph{Pedestrian Dead Reckoning.} Harle \citep{harle2013survey} surveys indoor PDR systems using smartphone IMUs. The Weinberg model \citep{weinberg2002using} estimates stride length from accelerometer variance as $s = k(a_{\max} - a_{\min})^{1/4}$. Jimenez et al.\ \citep{jimenez2009comparison} compare PDR algorithms and find that gyroscope heading integration with magnetometer correction can support short-duration tracking, but practical performance remains highly sensitive to phone placement, magnetic distortion, and calibration. Our implementation uses PDR as a bounded fallback, not as a long-duration standalone solution.

\paragraph{BLE Channel Sounding.} Bluetooth 6.0 (2024) introduces Channel Sounding, improving on RSSI-based ranging by combining phase-based and time-based measurements \citep{bluetooth2024cs}. Android 16 unifies UWB, BLE Channel Sounding, and WiFi NAN RTT under a single Ranging API \citep{android16ranging}. In practice, BLE Channel Sounding should be treated as an opportunistic decimeter-to-meter class input on commodity hardware, especially in cluttered indoor environments.

\paragraph{Decentralized trust.} Holochain \citep{holochain2024} provides agent-centric distributed hash tables where each participant maintains a sovereign data chain. Unlike blockchain, Holochain has no global consensus bottleneck and supports eventual consistency under arbitrary network partitions. In this work, Holochain is positioned as the trust and provenance fabric around the estimator, not as the estimator itself.

\section{System Architecture}

\subsection{Three-Layer Design}

The system consists of three independent layers:

\paragraph{Layer 1: Positioning Library.} A pure-Rust library with zero external dependencies (beyond \texttt{serde}), implementing the mathematical primitives: Gauss-Newton trilateration, a local Extended Kalman Filter, covariance-intersection helpers for decentralized peer fusion, Pedestrian Dead Reckoning, GDOP analysis, barometric altitude, multiple ranging models, and a sensor-agnostic measurement vocabulary. Body-agnostic via a \texttt{CelestialBody} trait with Earth (WGS-84), Moon, and Mars implementations.

\paragraph{Layer 2: Holochain Zomes.} Three zome pairs (anchor registry, ranging, position estimates) store anchor positions, range measurements, and fused position estimates on a distributed hash table. This layer is intended for admission control, provenance, delayed audit, and membership policy. It is not assumed to solve unknown-correlation fusion or to provide complete Sybil resistance through reputation weighting alone.

\paragraph{Layer 3: Mobile Bridge.} A Rust module compiled via JNI for Android (ARM64) that reads phone sensors (GPS, BLE, WiFi, UWB, accelerometer, gyroscope, barometer, magnetometer) and feeds them to the positioning library. It manages 5 positioning modes with automatic fallback and represents the first terrestrial adapter for a more general cooperative navigation core.

\subsection{Positioning Modes}

\begin{table}[h]
\centering
\caption{Positioning modes and their accuracy characteristics.}
\label{tab:modes}
\begin{tabular}{@{}llll@{}}
\toprule
\textbf{Mode} & \textbf{Condition} & \textbf{Accuracy} & \textbf{Duration} \\
\midrule
GPS Primary     & GNSS available          & 3--5 m       & Indefinite \\
Hybrid          & GNSS + peers            & 1--3 m       & Foreground use \\
Cooperative     & No GNSS, 3+ peers       & 0.3--3 m mixed radios$^\dagger$ & 1--3 min typical \\
Dead Reckoning  & No GNSS, no peers       & 10--50 m     & Tens of seconds to a few min \\
Offline         & No data, stationary     & Last known   & Indefinite \\
\bottomrule
\end{tabular}
\end{table}

\begin{flushleft}
\footnotesize{$^\dagger$BLE Channel Sounding performance is highly environment-dependent and should be modeled conservatively in cluttered indoor or NLOS conditions.}
\end{flushleft}

The mode transitions are automatic. When GPS is lost, the system saves the last known position as the PDR origin and begins integrating IMU data. If peer devices are discovered and admitted, the system transitions to Cooperative mode. When GPS returns, the dead-reckoned state is corrected. On stock mobile operating systems, the cooperative window should be assumed to depend strongly on foreground execution, permissions, and hardware support.

\section{Core Algorithms}

\subsection{Trilateration}

Given $N$ anchors at known positions $\mathbf{p}_i$ with range measurements $d_i$ and uncertainties $\sigma_i$, we minimize the weighted residual:

\begin{equation}
\hat{\mathbf{p}} = \arg\min_{\mathbf{p}} \sum_{i=1}^{N} \frac{w_i}{\sigma_i^2} \left( \|\mathbf{p} - \mathbf{p}_i\| - d_i \right)^2
\end{equation}

where $w_i \in [0,1]$ is a bounded quality prior for anchor $i$. We solve via iterative Gauss-Newton:

\begin{equation}
\Delta\mathbf{p} = (\mathbf{J}^T \mathbf{W} \mathbf{J})^{-1} \mathbf{J}^T \mathbf{W} \mathbf{r}
\end{equation}

where $\mathbf{J}$ is the Jacobian of residuals, $\mathbf{W} = \text{diag}(w_i/\sigma_i^2)$, and $\mathbf{r}$ is the residual vector. Convergence requires 3--5 iterations in practice. The covariance of the estimate is $\mathbf{C} = (\mathbf{J}^T \mathbf{W} \mathbf{J})^{-1}$.

\subsection{Bounded Pedestrian Dead Reckoning}

PDR estimates position from step counting and heading integration:

\begin{equation}
\mathbf{p}_{k+1} = \mathbf{p}_k + s_k \begin{bmatrix} \sin\psi_k \\ \cos\psi_k \end{bmatrix}
\end{equation}

where $s_k$ is the adaptive stride length and $\psi_k$ is the heading. Stride length follows the Weinberg model:

\begin{equation}
s_k = K (a_{\max,k} - a_{\min,k})^{1/4}, \quad K \approx 0.42
\end{equation}

Heading is integrated from the gyroscope rotation rate $\omega_z$:

\begin{equation}
\psi_{k+1} = \psi_k + \omega_{z,k} \cdot \Delta t
\end{equation}

When a magnetometer reading $\psi_m$ is available, heading is corrected via a low-gain complementary filter. In this prototype, PDR is explicitly treated as a short bridge between stronger external fixes rather than as a long-duration standalone estimator.

\subsection{Admission, Provenance, and Peer Fusion}

Each anchor's metadata may carry a quality prior derived from the Mycelix 8-dimensional sovereign profile stored on the Holochain DHT. For positioning-anchor trust specifically---where financial stake is irrelevant and the physics-relevant components dominate---we use a four-dimensional projection of the full profile:

\begin{equation}
w_i = \text{ep.\,integrity}_i^{0.25} \cdot \text{net.\,resilience}_i^{0.25} \cdot \text{civic}_i^{0.30} \cdot \text{competence}_i^{0.20}
\end{equation}

This projection is a subset of the full 8D credential defined in the governance papers; it drops dimensions (thermodynamic yield, economic velocity, stewardship, semantic resonance) that are not informative for positioning-anchor quality. This term can be used as an admission or measurement-quality prior, but it is not sufficient by itself to guarantee Sybil resistance or correlation-safe fusion. For peer estimates with unknown cross-correlation, conservative fusion should bound uncertainty explicitly. A standard covariance-intersection form is:

\begin{equation}
\Sigma_{\text{CI}}^{-1} = \omega \Sigma_1^{-1} + (1-\omega)\Sigma_2^{-1}
\end{equation}

\begin{equation}
\mu_{\text{CI}} = \Sigma_{\text{CI}} \left( \omega \Sigma_1^{-1}\mu_1 + (1-\omega)\Sigma_2^{-1}\mu_2 \right)
\end{equation}

In this architecture, Holochain carries provenance and admission policy, while the mathematical estimator remains responsible for conservative uncertainty management.

\subsection{Barometric Altitude}

Altitude is estimated from barometric pressure via the hypsometric formula:

\begin{equation}
h = 44330 \left(1 - \left(\frac{p}{p_0}\right)^{0.1903}\right) \text{ meters}
\end{equation}

This provides a useful vertical cue for floor transitions and local altitude changes, but its absolute accuracy remains weather-dependent.

\section{Results}

\subsection{Computational Performance}

All benchmarks were run on a single core in Rust release mode.

\begin{table}[h]
\centering
\caption{Computational performance benchmarks.}
\label{tab:benchmarks}
\begin{tabular}{@{}llll@{}}
\toprule
\textbf{Algorithm} & \textbf{Speed} & \textbf{Throughput} & \textbf{Error} \\
\midrule
3D Trilateration (6 anchors)  & 11.5 \textmu s/call & 87K/sec & 0.000 m \\
2D Trilateration (4 anchors)  & 1.7 \textmu s/call  & 588K/sec & 0.000 m \\
EKF (80 updates)              & 308 \textmu s/run   & 3.2K/sec & 0.000 m \\
PDR (500 steps)               & 0.3 \textmu s/step  & 3.3M/sec & 0.000 m$^*$ \\
GDOP computation              & $<$1 ns             & ---       & 1.528 \\
Ranging models                & $<$1 ns             & ---       & --- \\
Geodetic roundtrip            & 2.9 \textmu s       & 345K/sec & $<$0.01 m \\
\bottomrule
\end{tabular}
\begin{flushleft}
\footnotesize{$^*$PDR error is 0.000~m for perfect heading (square walk return-to-origin). In practice, useful PDR performance is strongly bounded by heading drift, phone placement, and magnetic distortion.}
\end{flushleft}
\end{table}

At a 20~Hz phone cycle rate, the total positioning computation requires $<$15~\textmu s per cycle (0.03\% of the 50~ms budget), leaving ample room for sensor fusion and rendering.

\subsection{Ranging Technology Comparison}

\begin{table}[h]
\centering
\caption{Supported ranging technologies and accuracy.}
\label{tab:ranging}
\begin{tabular}{@{}llll@{}}
\toprule
\textbf{Technology} & \textbf{Accuracy ($\sigma$)} & \textbf{Range} & \textbf{Hardware} \\
\midrule
UWB ToF              & 5--30 cm      & 50 m   & Pixel 6+, iPhone 11+ \\
BLE Channel Sounding & 20 cm--3 m    & 20 m   & BLE 6.0 (Android 16+) \\
WiFi RTT (802.11mc)  & 1--3 m        & 50 m   & Most modern routers \\
BLE RSSI             & 3--8 m        & 30 m   & All phones \\
LoRa ToA             & 50--200 m     & 10 km  & Meshtastic devices \\
Manual survey        & Instrument    & ---    & Any \\
Meshtastic hops      & 500--2000 m   & 50 km  & Meshtastic mesh \\
Optical ToF / laser  & mm--cm class  & mission-specific & space / robotics hardware \\
\bottomrule
\end{tabular}
\end{table}

BLE Channel Sounding is the most significant terrestrial addition because it improves materially on RSSI without requiring dedicated UWB silicon. It should nevertheless be modeled conservatively on commodity devices. UWB remains the stronger terrestrial anchor modality where available. For true cosmic deployments, optical and mission-grade RF timing are the more relevant analogues.

\subsection{Measurement Down-Weighting Under a Spoofed Prior}

In simulation, a single malicious anchor reporting ranges biased by 50~m was tested against 4 honest anchors. With uniform trust ($w_i = 1.0$ for all), the position estimate was displaced by 12.3~m. With the 4D-projected credential prior ($w_{\text{attacker}} = 0.01$, $w_{\text{honest}} = 0.90$), displacement was reduced to 0.4~m. This demonstrates that down-weighting helps under a spoofed-anchor prior; it does \emph{not} demonstrate complete Sybil resistance or solve unknown-correlation fusion in an open network.

\section{Discussion}

\subsection{The Bootstrap Problem}

A cooperative positioning system requires initial anchor positions. We address this with a ``GPS seed'' model: when GNSS is available, every participating device registers itself as a temporary anchor on the Holochain DHT with a geohash-indexed spatial key. When GNSS is denied, the cached anchor positions provide the reference frame. This creates a time-decaying positioning infrastructure that is strongest immediately after GNSS loss and degrades as anchor positions age.

\subsection{From Body-Agnostic Geometry to Cosmic Navigation}

The system is body-agnostic in the geometric sense: the \texttt{CelestialBody} trait abstracts geodetic-to-Cartesian transforms for any ellipsoidal or spherical body. That is necessary but not sufficient for interplanetary navigation.

A cosmic positioning system must replace the terrestrial phone sensor stack with mission-appropriate sensors: optical or laser inter-vehicle ranging, RF time transfer, star trackers for absolute attitude, inertial propagation, terrain-relative navigation for surface operations, and celestial timing or ephemeris-derived fixes for absolute reference. The value of the present work is that the mathematical core can already be expressed in sensor-agnostic terms.

\subsection{Limitations}

\begin{enumerate}
    \item PDR drift is unbounded without external correction and should be treated as a short bridge, not a long-duration indoor navigation solution.
    \item BLE Channel Sounding is available only on Bluetooth 6.0 hardware (2024+) and should be modeled conservatively under multipath and NLOS conditions.
    \item Stock mobile operating systems impose strict background-execution constraints on high-drain ranging technologies, limiting the plausibility of an always-on public mesh on commodity phones.
    \item A trust prior is not equivalent to a complete Sybil defense. Strong admission policy, identity proofs, and validation rules remain necessary in an open network.
    \item The current local EKF is appropriate for bounded local fusion but should not be treated as a complete solution for decentralized peer-state fusion with unknown cross-correlation.
    \item The system has not yet been validated with a real hardware campaign. All reported performance benchmarks are computational.
\end{enumerate}

\section{Conclusion}

We have implemented a decentralized cooperative positioning prototype with a real-time Rust computational core, explicit uncertainty models, and a decentralized provenance architecture. The system achieves lightweight runtime performance (11.5~\textmu s trilateration, 0.3~\textmu s PDR step) and is structured so that terrestrial and future non-terrestrial sensors can be expressed through the same estimator.

The most practically significant terrestrial result is that modern peer-to-peer ranging technologies can materially improve over RSSI-only positioning, provided their uncertainties are modeled honestly and bounded by conservative fusion logic. On smartphones, the current system is best described as a foreground-first bridge for localized GNSS-denied intervals rather than as a permanent infrastructure replacement.

When GPS fails, the system can degrade gracefully for a bounded interval: from satellite positioning to peer-to-peer ranging to dead reckoning, each mode handing off to the next as data sources are lost. The deeper contribution is architectural: by separating estimation, uncertainty, provenance, and sensor ingestion, the same kernel can evolve beyond phones.

The system is open-source (AGPL-3.0), implemented in Rust, and designed for deployment on Android today and more capable sensor stacks in the future. The same coordinate and fusion framework can be carried from Earth to the Moon, Mars, and deep-space missions once the terrestrial sensor assumptions are replaced with mission-grade ranging, attitude, and timing hardware.

\bibliographystyle{plainnat}

\begin{thebibliography}{20}

\bibitem[Android(2025)]{android16ranging}
Android Developers (2025).
\newblock Range between devices.
\newblock \url{https://developer.android.com/develop/connectivity/ranging}

\bibitem[Bluetooth SIG(2024)]{bluetooth2024cs}
Bluetooth SIG (2024).
\newblock Bluetooth Core Specification v6.0: Channel Sounding.
\newblock \url{https://www.bluetooth.com/specifications/specs/core60-html/}

\bibitem[Cannon et al.(2013)]{cannon2013extreme}
Cannon, P. et al. (2013).
\newblock Extreme space weather: impacts on engineered systems and infrastructure.
\newblock Royal Academy of Engineering.

\bibitem[DHS(2021)]{dhs2021pnt}
U.S. Department of Homeland Security (2021).
\newblock Positioning, Navigation, and Timing (PNT) Backup Requirements.

\bibitem[Harle(2013)]{harle2013survey}
Harle, R. (2013).
\newblock A Survey of Indoor Inertial Navigation Systems for Pedestrians.
\newblock \emph{IEEE Communications Surveys \& Tutorials}, 15(3):1281--1293.

\bibitem[Holochain(2024)]{holochain2024}
Harris-Braun, E. et al. (2024).
\newblock Holochain: A Framework for Distributed Applications.
\newblock \url{https://holochain.org}

\bibitem[ISPRS(2025)]{isprs2025uwb}
ISPRS Annals (2025).
\newblock UWB/Wi-Fi RTT Integration for Personal Mobility Applications.
\newblock ISPRS Annals X-G-2025.

\bibitem[Jimenez et al.(2009)]{jimenez2009comparison}
Jimenez, A.R. et al. (2009).
\newblock A Comparison of Pedestrian Dead-Reckoning Algorithms.
\newblock \emph{IEEE International Symposium on Intelligent Signal Processing}.

\bibitem[Weinberg(2002)]{weinberg2002using}
Weinberg, H. (2002).
\newblock Using the ADXL202 in Pedometer and Personal Navigation Applications.
\newblock Analog Devices Application Note AN-602.

\bibitem[Wymeersch et al.(2009)]{wymeersch2009cooperative}
Wymeersch, H., Lien, J., and Win, M.Z. (2009).
\newblock Cooperative Localization in Wireless Networks.
\newblock \emph{Proceedings of the IEEE}, 97(2):427--450.

\end{thebibliography}

\end{document}
